% General unary-frame and quotient-fibre criteria; no CCO geometry is assumed.
\section{Fresh real frames and hereditary unary amalgamation}
\label{sec:unary-fresh-frame-sap}

\subsection{Transferring orbit infinitude through real frames}

\begin{lemma}[Fresh frames and closed-base orbit infinitude]
\label{lem:foundational-real-frame-orbits}
Let $H$ be a many-sorted structure whose named functions are unary,
total on their specified source sorts, with no nullary generators.
Assume every finite set has finite unary closure. Let $G\le\Aut(H)$,
and let $R$ be a $G$-invariant union of real sorts generating all of $H$
under unary terms. Suppose that for every finite $D\subseteq R$ and
$u\notin\cl_H(D)$, the orbit $G_D\cdot u$ is infinite, where $G_D$
is the pointwise stabiliser. Then the following are equivalent:
\begin{enumerate}[label=(\roman*)]
\item For every finite unary-closed $A\subseteq B\subseteq H$ there is
  a finite $D\subseteq R$ with
  $A\subseteq\cl_H(D)$ and $\cl_H(D)\cap B=A$.
\item For every finite unary-closed $A\subseteq H$ and $u\notin A$,
  the orbit $G_A\cdot u$ is infinite.
\end{enumerate}
For an arbitrary finite hidden $T$, condition (ii) also says that the
points with finite $G_T$-orbit are exactly $\cl_H(T)$; every point of
that closure is fixed by $G_T$. No homogeneity or SAP is assumed.
\end{lemma}
\begin{proof}
Assume (i), fix $A,u$ as in (ii), and put
$B=\cl_H(A\cup\{u\})$. Choose a real frame $D$ from (i).
Then $u\notin\cl_H(D)$, so its $G_D$-orbit is infinite by the real-base
hypothesis. Every element of $G_D$ fixes all unary term values on $D$,
hence fixes $A$; thus $G_D\le G_A$, proving (ii).

Conversely assume (ii). Choose finitely many real generators $D_0$
with $A\subseteq\cl_H(D_0)$. Apply \cref{lem:finite-action-avoidance}
to $G_A$, with moving set $\cl_H(D_0)\setminus A$ and forbidden set
$B\setminus A$. Every moving orbit is infinite. For the resulting $g$,
put $D=gD_0$, which is still real because $G$ preserves $R$.
Since $g$ fixes $A$ and commutes with unary closure,
$A\subseteq\cl_H(D)=g\cl_H(D_0)$ and $\cl_H(D)\cap B=A$.
Finally $G_T=G_{\cl_H(T)}$, since unary term values are fixed whenever
all their generators are fixed. This proves the last assertion.
\end{proof}

\begin{remark}[Fresh real frames and hereditary hidden SAP]
\label{rem:typed-fresh-frame-iff-sap}
Suppose $H$ satisfies the structural hypotheses of
\cref{lem:unary-homogeneous-sap-criterion}, is generated by a union $R$
of named real sorts, and satisfies
$\operatorname{acl}_H(D)=\cl_H(D)$ for every finite real $D$.
Then the following are equivalent: fresh-frame separation as in
\cref{lem:foundational-real-frame-orbits}(i);
$\operatorname{acl}_H(T)=\cl_H(T)$ for every finite hidden $T$;
and SAP of the entire hereditary unary-closed age.
Indeed apply the preceding lemma with $G=\Aut(H)$, followed by
\cref{lem:unary-homogeneous-sap-criterion}. No CCO, quotient-lattice
geometry or functorial choice of representatives is required. For a
proper subgroup $G$, the preceding lemma retains the stronger statement
about its actual stabiliser orbits; SAP alone is not asserted to imply
that subgroup statement.
\end{remark}

\subsection{Finite avoidance within an actual quotient fibre}

\begin{lemma}[Quotient-fibre avoidance with named forced values]
\label{lem:foundational-quotient-fibre-avoidance}
Let $H$ have unary functions, total on their specified source sorts,
no nullary generators, and finite singleton unary closures. Let
$G\le\Aut(H)$ act transitively on a named real sort $M$, and let
$p_J:M\to Q$ be a named surjective equivariant quotient projection
with actual kernel $J>\Id_M$. For every well-sorted real-source unary
term $\tau:M\to H_s$, let $K_\tau$ be its actual kernel. Assume:
\begin{enumerate}[label=(\roman*)]
\item If $J\le K_\tau$, then $\tau=\sigma\circ p_J$ for an actual
  named unary term $\sigma:Q\to H_s$.
\item If $J\nleq K_\tau$, then $[J:J\cap K_\tau]=\infty$.
\end{enumerate}
For every $b\in Q$ and every finite forbidden set $Z\subseteq H$,
there is one $x\in p_J^{-1}(b)$ such that every unforced term
($J\nleq K_\tau$) satisfies $\tau(x)\notin Z$. Every forced value
is the prescribed $\sigma(b)\in\cl_H(\{b\})$. In particular
\[
             \cl_H(\{x\})\cap Z=\cl_H(\{b\})\cap Z,
             \qquad x\notin Z.
\]
The choices can be made successively for any finite list of such
quotient vertices, possibly on different transitive real sorts: at each
step the forbidden set may include any extra finite hidden set, all
earlier representatives and all earlier unforced term values. Thus the
real representatives can be pairwise distinct while all forced values
are retained.
\end{lemma}
\begin{proof}
Fix $x_0\in p_J^{-1}(b)$. There are only finitely many distinct unary
term functions from $M$: their values at $x_0$ lie in the finite
$\cl_H(\{x_0\})$, and equivariance with transitivity of $G$ makes
one such value determine the entire function on $M$.

The stabiliser $G_b$ of the quotient vertex is transitive on its actual
$J$-fibre. For an unforced term $\tau$, its orbit on $\tau(x_0)$ is
exactly $\tau(p_J^{-1}(b))$, whose cardinality is
$[J:J\cap K_\tau]=\infty$. Apply
\cref{lem:finite-action-avoidance} to the $G_b$-action on $H$, taking
as moving set all these finitely many unforced values $\tau(x_0)$ and
as forbidden set $Z$. One $g\in G_b$ avoids every forbidden target;
$x=gx_0$ is the required common representative. Forced terms take the
value $\sigma(b)$ throughout the fibre by (i).

Every element of $\cl_H(\{x\})$ is a real-source term value, so any
one in $Z$ is forced and belongs to $\cl_H(\{b\})$. Conversely
$b=p_J(x)$ gives $\cl_H(\{b\})\subseteq\cl_H(\{x\})$, proving the
intersection equality. The identity-real term has kernel $\Id_M$,
so it is unforced and gives $x\notin Z$. Successive use is legitimate
because every previously chosen representative has finite unary
closure; each enlarged forbidden set is still finite.
\end{proof}

\subsection{The relative-kernel criterion and its converse}

\begin{theorem}[Relative-kernel criterion for fresh frames and hereditary SAP]
\label{thm:typed-abstract-fresh-frame-sap}
Let $G$ act transitively on each of finitely many named original real
sorts $M_a$. Let $H$ be a countable homogeneous oligomorphic many-sorted
structure in a finite named language, with arbitrary relations and only
unary functions total on their specified source sorts, no nullary
generators, and finite unary closures of finite sets. Assume
$G\le\Aut(H)$ preserves all named sorts, predicates and functions.
The singleton closures on the finitely many transitive real sorts are
then uniformly bounded; no uniform bound independent of the number
of generators is required. We use orbit-theoretic $\acl_H$ as in
\cref{lem:unary-homogeneous-sap-criterion}.

Every nonreal hidden vertex is assumed to have a representation
$p_J(x)=[x]_J$ for some named surjective equivariant quotient projection
$p_J:M_a\to Q_{a,J}$, where $J>\Id_a$ is an actual $G$-invariant
equivalence. For every such $J$ and every named real-source unary
term $\tau:M_a\to H_s$ with kernel $K_\tau$, assume the exact dichotomy:
\begin{enumerate}[label=(\roman*)]
\item If $J\le K_\tau$, then $\tau=\sigma\circ p_J$ for an actual
  named unary term $\sigma:Q_{a,J}\to H_s$.
\item If $J\nleq K_\tau$, then $[J:J\cap K_\tau]=\infty$.
\end{enumerate}
Finally assume that for every finite original real $D$ and every
$u\notin\cl_H(D)$, the orbit $\Aut(H/D)\cdot u$ is infinite.

Then for every finite induced unary-closed $A\subseteq B\subseteq H$
there is a finite original real $D$ such that
\[
       A\subseteq\cl_H(D),\qquad \cl_H(D)\cap B=A.
\tag{AFF}\label{eq:typed-abstract-fresh-frame}
\]
Consequently $\acl_H(T)=\operatorname{dcl}_H(T)=\cl_H(T)$ for every
finite hidden $T$, and the entire hereditary age of finite unary-closed
substructures of $H$ has SAP.

This criterion assumes no original CCO, triangle geometry, lattice
distributivity or old SAP. In the terminal application those conditions
supply the named quotient signature, kernel dichotomy and real-base
orbit condition separately. Kernel inclusion alone cannot replace the
literal named factorisation in (i).
\end{theorem}
\begin{proof}
Start with the real points $D_0$ of $A$. For each nonreal $b_i\in A$,
choose a named representation $b_i=p_{J_i}(x_i^0)$. Apply
\cref{lem:foundational-quotient-fibre-avoidance} successively, with a
forbidden set containing $B$ and the previously chosen real
representatives. It gives $x_i$ in the prescribed fibre with every
unforced term value outside $B$, while every forced value belongs to
$\cl_H(\{b_i\})\subseteq A$.

Put $D=D_0\cup\{x_i\}$. It generates all of $A$, and its closure is
the union of the singleton closures of its real generators. The
$D_0$ closures lie in $A$; each new singleton closure meets $B$ only
in $A$. This proves (AFF), including the empty-base case.

Apply \cref{lem:foundational-real-frame-orbits} with the full group
$\Aut(H)$ and the union of the named real sorts. The projection
representations give real generation, and the theorem's real-base
hypothesis supplies the required orbits. Thus every vertex outside
any finite closed $A$ has infinite orbit over $A$. The final conclusions
follow from \cref{lem:unary-homogeneous-sap-criterion}.
\end{proof}

\begin{corollary}[Exact abstract relative-kernel characterisation of unary hereditary SAP]
\label{cor:typed-abstract-relative-kernel-sap-iff}
Keep the structural hypotheses on $G$, the original transitive real
sorts, $H$, its finite unary closures and named surjective quotient
representations from \cref{thm:typed-abstract-fresh-frame-sap}, without
assuming its kernel dichotomy or real-base orbit property. Then the
following are equivalent:
\begin{enumerate}[label=(\roman*)]
\item The entire hereditary age of finite induced unary-closed
  substructures of $H$ has SAP.
\item For every finite original real $D$, $\acl_H(D)=\cl_H(D)$;
  and for every named quotient $p_J:M_a\to Q_{a,J}$ and every actual
  named real-source unary term $\tau:M_a\to H_s$ with kernel $K_\tau$,
\[
\begin{array}{ll}
 J\le K_\tau:
  &\exists\text{ named unary term }\sigma:
        \tau=\sigma\circ p_J,\\[2pt]
 J\nleq K_\tau:
  &[J:J\cap K_\tau]=\infty.
\end{array}
\tag{RK}\label{eq:typed-abstract-kernel-iff}
\]
\end{enumerate}
The forward direction needs no separate original CCO, SAP, exact-lattice
or real-base orbit input. The two failure mechanisms are local to one
quotient fibre: a missing named factorisation for an existing forced
real-source term gives a definable off-closure value, while a finite
proper relative index for an existing term kernel gives an algebraic
off-closure value. The quantification over actual named terms is essential.
\end{corollary}
\begin{proof}
Condition (ii) gives precisely the hypotheses of
\cref{thm:typed-abstract-fresh-frame-sap}: finite $D$ and its unary
closure have the same pointwise stabiliser. Thus (ii) implies (i).

Conversely assume (i). By
\cref{lem:unary-homogeneous-sap-criterion}, $\acl_H(T)=\cl_H(T)$
for every finite hidden $T$, in particular every finite real $D$.
Fix $p_J,\tau$, choose $x\in M_a$, and put $b=p_J(x)$, $v=\tau(x)$.

If $J\le K_\tau$, the term $\tau$ is constant on the entire fibre of
$b$. Every $g\in\Aut(H/b)$ preserves that fibre, so
$gv=\tau(gx)=v$. Hence $v\in\acl_H(\{b\})=\cl_H(\{b\})$,
and $v=\sigma(b)$ for an actual named unary term $\sigma$.
Equivariance and transitivity of $G$ on $M_a$ extend this equality at
one $x$ to $\tau=\sigma\circ p_J$ on all of $M_a$.

If $J\nleq K_\tau$ and $[J:J\cap K_\tau]=r<\infty$, then $r\ge2$.
The set
\[
             V_b=\{\tau(y):y\in M_a,\ p_J(y)=b\}
\]
has exactly $r$ elements and is invariant under $\Aut(H/b)$, so
$v\in\acl_H(\{b\})$. But $G_b$ is transitive on the fibre and hence
on $V_b$; since $r\ge2$, some element of $G_b$ moves $v$. Every named
unary term value of $b$ is fixed by $G_b$, so
$v\notin\cl_H(\{b\})$, a contradiction. Therefore the unforced
relative index is infinite, proving (ii).
\end{proof}
